% ch4_policy.tex — 4장 정책 제언
\section{정책 제언}
\label{sec:policy}

\noindent{\sffamily\small\color{muted}\textbf{CHAPTER 4}}

\begin{quoteblock}
관측\textperiodcentered{}계수\textperiodcentered{}텍스트의 3중 신호를 결합한 장소민감적 정책 처방을 제시하고, 본 프레임워크의 외적 타당성과 후속 연구를 위한 조건을 정리한다.
\end{quoteblock}

\figfull{image21.png}{장소민감적 정책 제언 4대 축}{fig:0-5}

본 장에서는 3장의 분석 결과에 기반하여 서울의 디지털 격차 해소를 위한 정책 방향을 제시한다. 분석에 활용된 데이터는 2023--2024년에 수집된 것으로, 지능정보화 기본법 체제 하의 정책 환경을 반영한다. 2026년 1월부로 디지털포용법이 시행되어, 기존의 정보격차 해소 중심 접근에서 디지털포용이라는 확장된 개념으로 정책 프레임이 전환되었다. 본 장의 제언은 기존 체제 하에서 확인된 격차와 취약성을 실증적 근거로 제시하고, 새로운 법률 체계 내에서의 정책적 활용 가능성을 모색하는 구조를 취한다.

\subsection{정책 제언의 이론적 프레임}
\label{subsec:4-1}

서울의 디지털 격차에 대한 정책 접근은 \textbf{people-based}와 \textbf{place-based}의 두 축으로 구분할 수 있으며, 3장의 분석 결과는 전자가 정책의 주축으로 적합함을 시사한다. 3.2절에서 개인 수준 변수만을 투입한 Model 1이 자치구 간 분산의 83.4\%를 설명하였는데, 이는 자치구 간 디지털 활용 격차의 대부분이 지역 고유의 특성보다 해당 지역에 거주하는 주민의 인구사회학적 구성 차이에 기인함을 의미한다. 따라서 정책의 일차적 개입 경로는 저학력 고령층, 독거 고령자 등 복합적 취약 요인을 가진 집단의 디지털 역량을 직접 강화하는 people-based 접근에 놓아야 한다.

다만 여기서의 place-based 접근은 자치구 단위의 일률적 자원 배분과 구별되어야 한다. ICC 0.49\%는 자치구 간 분산이 디지털 활용 격차의 주된 원천이 아님을 나타내며, 이는 지역 단위 배분 중심의 정책이 비효율적일 수 있음을 시사한다. 본 프로젝트에서 제안하는 것은 장소민감형(place-sensitive) 접근, 즉 people-based 정책의 효과를 공간적 조건 정비를 통해 증폭시키는 전략이다. 교차수준 상호작용에서 확인된 바와 같이, 공공 WiFi 밀도가 고령자에게(+3.03점), 연결성 수준이 저학력자에게(+10.94점) 불균형적 효과를 나타냈다는 것은, 취약계층 밀집 지역의 인프라 수준 제고가 people-based 정책의 핵심적인 승수로 기능할 수 있음을 의미한다.

이러한 이중 접근은 국제적으로도 통용되는 정책 프레임과 일관된다. \citet{OECD2025}는 지역 간 격차 확대(uneven divergence) 구조에서 장소 맞춤형 정책과 사람 중심 이전 정책이 상호 배타적이 아닌 보완적 관계임을 강조하며, \citet{ITU2024}의 UMC 정책 어젠다 역시 타깃 그룹 중심 정책과 국가 디지털 전략\textperiodcentered{}로드맵 수립을 동시에 권고한다. EU Digital Decade 2030은 디지털 기술\textperiodcentered{}인프라\textperiodcentered{}비즈니스\textperiodcentered{}공공서비스의 4개 축에 대해 회원국별 국가 로드맵을 수립하고, 연례 모니터링을 통해 진척을 점검하는 순환 구조를 운용하고 있다~\citep{EuropeanCommission2025}. 이러한 측정 $\rightarrow$ 모니터링 $\rightarrow$ 정책 조정의 구조는 본 프로젝트의 분석 절차와 일부 부합한다.

\subsection{People-based 정책 제언}
\label{subsec:4-2}

\subsubsection*{(1) 독거 고령자}

독거 고령자의 디지털 활용 예측 격차는 약 $-10.87$점으로, 전체 척도의 13.6\%에 해당한다. 다인가구 고령자가 가구원으로부터 비공식적 디지털 지원을 받을 수 있는 반면, 독거 고령자에게는 이 경로가 부재하다. 3.3절의 텍스트 분석에서 이 격차의 일상적 발현이 확인되었다. 송파구 문정동의 한 IT 기획자가 ``이 글을 보시는 분은 어느 정도 충분히 잘 쓰시고 계신다는 말이겠죠, 자녀분들이 없으시거나 아님 떨어져 살아서 묻기 힘들거나''라며 중장년\textperiodcentered{}노년층 대상 앱 사용법 교육을 제안한 사례는, 가구 내 디지털 지원망의 부재가 일상적 디지털 과업의 지속적 어려움으로 이어지는 구조를 가시화한다.

기존 디지털배움터는 집합 교육 중심으로, 이동이 제한적인 독거 고령자에 대한 물리적 도달에 구조적 한계가 있었다. 2026년 3월부터 AI디지털배움터로 전환되면서 찾아가는 파견교육이 확대될 예정이나(과학기술정보통신부, 2025.12.24), 독거 고령자에 특화된 프로그램 설계는 아직 구체화되지 않았다. 영국 Age UK의 Digital Champion Programme이 자원봉사 기반 1:1 가정방문과 기기 대여를 병행하여 고령자의 디지털 자신감 향상을 보고한 바 있다~\citep{AgeUK2023}.

이를 종합하면, 독거 고령자에 대한 정책은 집합 교육의 접근성 확대만으로는 불충분하며, 사회적 비계를 대체하는 모델이 필요하다. 구체적으로, 동 단위 디지털 동반자(Digital Companion) 양성을 통한 1:1 방문 지원 체계, 독거 고령자 특화 코호트 운영, 기기 대여 병행을 제언한다.

\subsubsection*{(2) 저학력 고령층}

M3에서 연령 $\times$ 저학력 상호작용($-0.28$, $p < 0.001$)은 저학력 집단의 연령별 디지털 활용 감소 기울기가 대졸 기준 대비 약 2배 가파름을 보여준다. 20년의 연령 범위에 걸쳐 이 차이는 약 5.5점의 추가 격차로 누적되며, 교육 집단 간 디지털 격차가 연령 증가에 따라 확대되는 경향을 시사한다.

3.3절의 텍스트 분석은 이 교육 기울기가 일상에서 어떻게 발현되는지를 보여준다. `정보부재'가 가장 빈번한 결여 유형(63\%)으로 나타난 것은, 교육수준이 낮을수록 디지털 서비스의 존재 자체를 인지하지 못하는 격차가 심화됨을 시사한다. 편의점 무인 인쇄 서비스를 모르는 사례, ISP 고객센터 전화번호를 모르는 사례, 정부24의 공문서 출력 기능을 모르는 사례 등은 모두 `서비스가 존재하지만 이용자에게 보이지 않는' 상태를 가리킨다. 이러한 정보부재는 교육수준과 강하게 연관될 수 있는데, 검색 역량(informational skill)이 학력 수준에 비례하기 때문이다~\citep{Hargittai2010}.

기존 디지털배움터가 교육수준별로 차등화된 경로를 체계적으로 제공했는지는 불분명하다. 싱가포르 IMDA의 Seniors Go Digital은 기초--중급--고급의 3단계 경로를 설정하고, 각 단계에 사이버 보안 교육을 내장하며, 저소득 고령자 대상 기기\textperiodcentered{}요금 보조를 병행하는 통합 모델을 운용한다~\citep{IMDA2020, CSASingapore2023}. 특히 보안 교육은 3.1절에서 Safety 차원의 취약성이 확인된 서울 맥락에서 참고할 만하다.

제언으로는, 교육수준에 따른 단계별 경로의 체계적 차등화, 각 단계에 보안 교육을 필수로 포함하는 통합 설계, 저소득 저학력 고령층을 위한 기기\textperiodcentered{}요금 지원 연계를 제시한다.

\subsubsection*{(3) 정보부재 해소를 위한 디지털 서비스 발견가능성 강화}

3.3절의 텍스트 분석이 제기하는 정책적 과제 중 하나는, 기존의 디지털 역량 교육이 `사용법(how to use)'에 집중하는 반면, `존재 인지(awareness of existence)'의 격차는 정책적으로 충분히 다루어지지 않았다는 점이다. 분석 대상 100건 중 63건에서 추론된 정보부재는, 서비스를 사용할 역량이 아니라 서비스의 존재 자체를 모르는 상태를 가리킨다.

예컨대, 관악구 봉천동의 주민이 ``법원서류 등은 일반 프린트카페에서 인쇄가 안 되더라구요, 공문서 프린트 가능한 곳 아시는 분''이라고 요청한 사례에서, 정부24의 공문서 출력 기능이나 법원 무인발급기의 존재 자체가 이용자에게 가시적이지 않았다. 은평구 신사동의 ``정녕 없는 건지요''라는 표현은 온라인 검색을 이미 시도했으나 결과가 불충분했을 가능성을 시사하며, 인프라의 존재와 인프라의 `발견가능성(discoverability)' 사이의 간극을 조명한다.

이러한 정보부재는 단순한 홍보 강화로 해결되기 어렵다. 정보를 수동적으로 제공하는 방식은 이미 정보 탐색 역량을 갖춘 이용자에게만 도달하기 때문이다. 따라서 정책의 설계 원칙으로서, 디지털 서비스의 발견가능성을 서비스 설계 단계에서부터 내장하는 접근이 필요하다. 구체적으로, 공공 서비스 앱(서울온 등)에서 위치 기반으로 인근 공공 디지털 접근점(무인 인쇄, 도서관 복합기, 디지털배움터)을 안내하는 기능, 주민센터 방문 민원 시 관련 디지털 서비스를 안내하는 오프라인 연계, 관할 동사무소의 문자 알림 시 관련 디지털 경로를 병기하는 방안 등을 제언한다.

\subsubsection*{(4) 교육 기울기에 대한 구조적 대응}

교육수준은 디지털 활용의 가장 강력한 예측 변수로, 이 격차는 단기 프로그램만으로 해소되기 어려운 구조적 성격을 갖는다. 따라서 정책의 시간 지평을 구분할 필요가 있다. 단기적으로는 상기 타겟형 프로그램이 현재 격차의 완화에 기여하고, 중장기적으로는 기초 교육 단계에서부터 디지털 역량을 교육 과정에 체계적으로 포함하는 접근이 향후 세대의 격차 누적을 예방하는 구조적 경로가 된다. 디지털포용법 제8조의 기본계획이 요구하는 `디지털역량 함양 활동에 관한 사항'과 `디지털포용을 위한 교육\textperiodcentered{}상담 및 홍보'는 바로 이러한 단기\textperiodcentered{}중장기 정책을 동시에 포괄하는 제도적 기반을 제공한다.

\subsection{Place-based 정책 제언}
\label{subsec:4-3}

\subsubsection*{(1) 디지털 사막 자치구의 차원별 맞춤 개입}

3.1절에서 종합지수 하위 5개 자치구(중랑구, 도봉구, 강북구, 구로구, 노원구)를 `디지털 사막'으로 분류하였다. 이들의 핵심 특징은 Connectivity 차원에서는 도시 평균 수준 이상을 기록하면서도(중랑구 0.630, 강북구 0.562), Devices(중랑구 0.052)와 Safety(도봉구 0.019) 차원에서 극단적 결핍을 보인다는 점이다. 이는 물리적 네트워크 인프라의 부재가 아닌, 이를 활용하는 인적\textperiodcentered{}경제적 역량의 부족이 주요 격차 요인임을 시사한다.

차원별 분해(그림 5)에서 각 자치구의 취약 차원이 상이한 것으로 확인되었으므로, 일률적 지원보다는 차원별 맞춤 개입이 보다 효과적일 수 있다. 예컨대, Devices 차원이 극단적으로 낮은 중랑구에는 기기 보급\textperiodcentered{}대여 프로그램이, Safety 차원이 사실상 0에 가까운 도봉구에는 보안 교육\textperiodcentered{}인식 제고 프로그램이 우선적으로 필요하다. Affordability 차원이 취약한 자치구에는 요금 보조\textperiodcentered{}저가 요금제 확대가 보다 직접적인 개입 경로가 된다.

3.3절의 텍스트 분석은 차원별 맞춤 개입의 구체적 방향을 제시한다. Safety 차원이 극단적으로 낮은 도봉구의 경우, 텍스트 분석에서 온라인 중고거래 피해 사례가 안전 결여의 주요 발현 형태로 나타났다. 중고 전자기기 거래에서 상태 검증 역량의 부재와 피해 후 공식 구제 경로(소비자상담센터, 경찰 신고 등)에 대한 정보부재가 결합된 패턴은, 보안 교육이 단순한 비밀번호 관리를 넘어 디지털 거래 환경에서의 소비자 보호 역량까지 포괄해야 함을 시사한다.

M2에서 자치구 인프라 변수의 주효과가 전체적으로 유의하지 않았음($p = 0.237$)을 고려하면, place-based 개입이 단독으로 격차를 해소하기에는 한계가 있다. 이는 차원별 맞춤 개입이 4.1.1의 people-based 정책과 결합될 때 --- 예컨대, 기기 보급과 함께 수준별 교육이 병행될 때 --- 보다 실질적인 효과를 기대할 수 있음을 의미한다.

현행 정책 체계에서 자치구별 특성을 반영한 차등적 자원 배분이 이루어지고 있는지에 대한 검토가 필요하다. 영국의 Digital Inclusion Innovation Fund(\pounds11.7M, 2025)는 80개 지역 프로젝트에 커뮤니티 주도 방식으로 자금을 배분하되, 각 지역의 필요(노인, 저소득층, 이주민 등)에 따라 사업 내용을 차별화하는 모델을 채택하였다~\citep{UKDSIT2025}. 이는 지역 맥락에 맞는 차등적 개입의 실현 가능성을 보여주는 참고 사례이다.

\subsubsection*{(2) 공공 디지털 접근점의 재설계}

3.3절의 텍스트 분석에서 공공 접근점 결여가 32\%의 게시글에서 추론된 것은, 서울의 공공 디지털 인프라가 양적으로 부족한 것이 아니라 이용자의 과업 맥락에 부합하지 않는 방식으로 설계\textperiodcentered{}배치되어 있을 가능성을 제기한다. 관악구 신림동의 자취생이 ``문서 스캔이 안 되는데 되는 PC방 추천''을 요청한 사례나, 관악구 봉천동의 주민이 ``사당역 근처 인터넷, 프린터 쓸 수 있는 곳 어디 없나용''이라 질문한 사례에서, 편의점 무인 인쇄\textperiodcentered{}도서관 복합기 등 기존 공공 접근점의 존재가 이용자에게 인지되지 않았다.

이러한 결과는 공공 디지털 접근점 정책이 `설치'에서 `발견가능성과 접근성'으로 전환되어야 함을 시사한다. 구체적으로, 자치구별 공공 디지털 접근점(무인 인쇄, 스캔, 공공 WiFi, 디지털배움터)의 통합 안내 체계 구축, 네이버 지도\textperiodcentered{}카카오맵 등 민간 플랫폼과의 데이터 연계를 통한 검색 가능성 강화, 심야\textperiodcentered{}주말 등 공공시설 운영시간 외 시간대의 접근점 확보(24시간 무인 인쇄 부스 등)를 제언한다.

\subsubsection*{(3) 서울 북부 권역 연합 거버넌스}

UMC 종합지수의 공간적 분포(그림 1)와 LISA 클러스터 분석(그림 5)은 디지털 사막 자치구들이 서울 북부에 통계적으로 유의하게 군집해 있음을 보여준다. Low-Low 클러스터에 해당하는 도봉구, 강북구, 노원구, 중랑구 등이 물리적으로 인접해 있다는 사실은, 이 문제가 개별 자치구의 문제가 아닌 권역 수준의 구조적 패턴임을 시사한다.

이러한 공간적 군집은 권역 단위 협력 거버넌스를 제언하는 세 가지 근거를 제공한다. 첫째, LISA 분석에서 공간적 자기상관이 유의한 것으로 확인되어, 단순한 시각적 패턴이 아닌 구조적 현상임이 통계적으로 뒷받침된다. 둘째, 개별 자치구가 각자 디지털 교육 프로그램이나 인프라를 독립적으로 구축할 경우 자원이 분산되지만, 인접 자치구가 공동으로 디지털 허브를 운영하거나 교육 프로그램을 공유하면 규모의 경제를 실현할 수 있다. 셋째, 3.1절에서 확인된 바와 같이 인접 자치구 간에도 차원별 강점과 약점이 상이하여(예: 중랑구의 Connectivity 0.630 강점 vs Devices 0.052 약점), 자원의 교차 활용이 가능하다.

구체적으로, 북부 디지털 사막 권역(도봉\textperiodcentered{}강북\textperiodcentered{}노원\textperiodcentered{}중랑 등)의 연합 거버넌스 구축을 제언한다. 이는 공동 디지털 교육 허브의 운영, 차원별 강점-약점의 교차 활용, 권역 단위 예산 확보 및 사업 기획을 포함한다. 이러한 권역 단위 협력 체계의 제도적 경로는 4.1.3에서 디지털포용법의 기본계획 체계와 연결하여 논의한다.

\begin{keymsg}[title=정책 시사점 \textperiodcentered{} 장소민감적 처방의 4대 축]
디지털 사막 우선 개입 $\rightarrow$ 차원별 차등 처방 $\rightarrow$ 사람\textperiodcentered{}장소 통합 접근 $\rightarrow$ 정량+텍스트 모니터링. 이 4대 축은 3장의 측정\textperiodcentered{}HLM\textperiodcentered{}텍스트 결과를 하나의 정책 프레임으로 통합한다.
\end{keymsg}

\subsection{통합적 정책 제언}
\label{subsec:4-4}

앞선 분석 결과와 정책 제언은, people-based와 place-based 접근이 결합될 때 가장 실질적인 효과를 기대할 수 있음을 시사한다. 공간적 타겟팅과 집단적 타겟팅의 결합 --- 디지털 사막 자치구에 독거 고령자 찾아가는 지원, 수준별 교육 프로그램 등 people-based 정책을 우선 배치하는 방식 --- 이 구체적인 통합 경로가 된다. 3.3절의 텍스트 분석은 이 통합의 구체적 양상을 보여준다. 복합 결여(88\%)가 보편적이라는 발견은, 단일 축의 정책 개입---기기만 보급, 교육만 제공, 인프라만 확충---이 결여의 중첩 구조를 해소하기에 불충분함을 시사한다. 예컨대, 기기를 보급하더라도 해당 기기로 접근 가능한 서비스의 존재를 모르면(정보부재), 또는 서비스의 디지털 인터페이스를 조작할 역량이 없으면(역량부재), 기기 보급의 효과는 제한적이다.

디지털포용법 제8조는 정부가 3년마다 디지털포용 기본계획을 수립하도록 규정하며, 기본계획에는 `디지털취약계층의 접근성 개선 및 사회 참여 지원', `디지털포용의 실태와 정책성과 분석', `디지털역량 함양 활동' 등이 포함되어야 한다. 동법 제9조는 관계 중앙행정기관 및 지방자치단체가 매년 시행계획을 수립\textperiodcentered{}시행하도록 하고 있다. 본 프로젝트에서 수행한 다층 분석 --- 구 간\textperiodcentered{}구 내 분산 분해, 취약집단 식별, 차원별 프로파일링 --- 의 결과는 이 기본계획 수립 과정에서 실증적 근거로 활용될 수 있다. 특히 3.3절의 텍스트 분석이 도출한 결여 유형별 분포(정보부재 63\%, 제도적 경로 결여 61\%, 역량 부재 45\%)는 기본계획의 우선순위 설정에 참고할 수 있는 경험적 근거를 제공한다.

4.1.2에서 제안한 북부 디지털 사막 권역 연합 거버넌스의 제도화 역시 이 체계 안에서 모색될 수 있다. 디지털포용법은 광역 및 기초 지방자치단체를 시행계획 수립\textperiodcentered{}시행의 주체로 규정하고 있으나, 인접 기초자치단체 간 권역 단위 협력 체계에 대한 별도의 규정은 부재한다. 기본계획 수립 과정에서 이러한 권역 단위 협력의 제도적 경로 --- 예컨대 자치구 간 공동 사업 기획을 위한 협의체 구성 등 --- 가 검토될 필요가 있다.

마지막으로, 1.1에서 논의한 절차적 일반화 가능성의 관점에서, 본 프로젝트의 분석 체계 --- 디지털 격차 진단(종합지수 구축 및 차원별 분해) $\rightarrow$ 다층 분석(개인\textperiodcentered{}지역 기여 분리 및 취약집단 식별) $\rightarrow$ 질적 맥락 추출(LLM 기반 결여 유형 분류) $\rightarrow$ 타겟 정책 설계(people-based + place-based 통합) --- 는 서울 특정의 처방을 넘어, 유사한 데이터 구조를 갖는 다른 도시에도 적용 가능한 정책 도출 절차로서 의의를 갖는다.
